% Ad Familiares 2.19 --- headnote
% ad-familiares-02-19, 22 June 50 BC, camp in Cilicia

\textit{Cicero to C. Caelius Caldus, his quaestor for the Cilician
year, written in camp in Cilicia on the tenth day before the Kalends
of Quintilis (22 June) 50 BC (the manuscript dateline: \textit{Scr.
in castris in Cilicia a. d. x K. Quint. a. 704 (50)}). The formal
opening sets the tone --- both correspondents are named with full
filiation, \textit{M. Tullius M. f. M. n. Cicero imp.} writing to
\textit{C. Coelius L. f. C. n. Caldus q.} --- the salutation of
proconsul to quaestor on first acquaintance.}

\textit{Caldus had drawn Cicero by lot, and Cicero had been pleased
with the assignment; but the young man had been slow to set out, and
no word had reached the province about his coming until a letter
arrived in camp, dated the tenth day before the Kalends of Quintilis
but otherwise without an indication of place, day of despatch, or
expected arrival. Cicero now writes back by his own orderlies and
lictors to press Caldus to hurry: Cicero's year is nearly out and he
fears he will be gone before Caldus reaches him. The letter is
diplomatic in the way an introductory letter must be. Recommendations
from Curius and from C. Vergilius have arrived, and weigh with Cicero;
but Caldus's own letter weighs more, and ``no quaestor more welcome
could have fallen to me.'' The closing pledge --- that whatever
distinctions Cicero confers shall mark the honour of Caldus and of
his ancestors --- restates the courteous bargain of the proconsul's
year, with one private note of warning embedded in it: Cicero will
discharge that obligation \textit{the more easily} if Caldus reaches
him in time.}

\textit{The text at the head of section 2 carries a crux
(\textit{quaecumque a me ornamenta \dag ad te proficiscentur}); the
sense of the period is clear enough.}
